Chapter~\ref{ch:ir} introduced the Cranfield paradigm: finding the
document that best matches what you're looking for.  It also showed
how ``good old-fashioned'' \abr{ai} systems use rules and logic to
answer questions about specific domains.
%
Much of question answering was built on information retrieval: given a
question, you want to find the Wikipedia page or other information
source that answers the question.

However, in the 2010s, retrieval was reinvented by new neural
techniques that revolutionized not just question answering but
computer science more broadly.
%
You may have heard of these models like \camelabr{Chat}{gpt} or \abr{bert}.
%
The ``T'' in both of these models stands for ``transformer'', the core technology that we'll discuss in this chapter and the next.
%
Because this is an important distinction, I'm going to divide this into two parts: methods that primarily use an ``encoder'' architecture and those that focus on the ``decoder''.  
%
But what are these models encoding to and from?  
%
Vectors.

We've already talked about vectors earlier in the book: Chapter~\ref{ch:ir} used vectors to help find the right document to answer a question. 
%
But those vectors had dimensions that were mostly empty: there was a dimension for every word, and only the important words were non-zero.  
%
This so-called ``sparse'' encoding was easier for humans to understand, but didn't work as well as the ``dense'' vectors we'll talk about in this chapter.
%
Here, the dimensions are \emph{learned}: the dimensions are the ones that help to divide up the dataset to help the retrieval or other mechanisms find the right documents.

The vectors build using the methods described in this chapter were the precursors of the world-changing models
we discuss in the next chapter.
% 
Those methods were necessary because of the pressures outlined in this chapter: you needed vectors to not just show what document had the answer but rather to explain \emph{how} to answer question.

% TODO(jbg): Outline the entire chapter

\section{But What About Non-Text?}
\label{sec:methods:multimodal}

Thus far, all of the questions we've talked about in this chapter have been text.  But most of the methods naturally extend to other modalities.  As we discussed in Section~\ref{sec:formats:multimodal}, multimodal encoders can transform multiple languages, images, and audio into a consistent representation.  Those become the same kind of input to the models that predict the next word.  Those models can answer questions directly or be ``fine-tuned'' to mimic better answers or fine-tuned to produce chain of thought analysis before their final answer.


